\subsection{Römer 6}
\subsubsection{Text}
\textbf{1} Was \vb{wollen}\annot{Präsens, Indikativ, Aktiv} wir \konj{nun}\annot{Gedankenschritt} \vb{sagen}\annot{Infinitiv}? 

    \spa{1}\vb{Sollen}\annot{Präsens, Indikativ} wir in der Sünde \vb{verharren}\annot{Infinitiv}, 

        \spa{2}\konj{damit}\annot{Ziel} das Mass der Gnade \vb{voll werde}\annot{Präsens}?

    \spa{0}\textbf{2} Das \vb{sei}\annot{Konjunktiv} ferne!\annot{Anwort auf die 2 rethorische Fragen}

    \spa{1}Wie \vb{sollen}\annot{Präsens, Indikativ} wir, 

        \spa{2}die wir der Sünde \vb{gestorben sind}\annot{Perfekt, Indikativ}, 

        \spa{2}\konj{noch}\annot{Grund} in ihr \vb{leben}\annot{Präsens, Infinitv}?

    \spa{1}\textbf{3} \konj{Oder}\annot{Besondere Betonung} \vb{wisst}\annot{Präsens, 2. Person Indikativ} ihr nicht, 

        \spa{2}\konj{dass}\annot{Folge aus den Fragen} wir alle, 

        \spa{2}die wir in \person{Christus Jesus} \verbP{hinein getauft sind}\annot{Perfekt Indikativ}, 

        \spa{2}in seinen Tod \verbP{getauft sind}\annot{Perfekt Indikativ}?

    \spa{1}\textbf{4} Wir \vb{sind}\annot{Präsens, Indikativ} \konj{also}\annot{Ursache} mit ihm \verbP{begraben worden}\annot{Perfekt, Indikativ} 

    \spa{1}\konj{durch}\annot{Art} die Taufe in den Tod, 

        \spa{2}\konj{damit}\annot{Ziel}

            \spa{3}\konj{gleichwie}\annot{Ergebnis} \person{Christus} 

            \spa{3}\konj{durch}\annot{Mittel} die Herrlichkeit des \person{Vaters} aus den Toten \verbP{auferweckt worden ist}\annot{Perfekt, Indikativ}, 

        \spa{2}\konj{so auch}\annot{Art und Weise} wir in einem neuen Leben \vb{wandeln}\annot{Präsens, Indikativ, Infinitiv}.
        
    \spa{1}\textbf{5}  \konj{Denn}\annot{Erklärung} wenn wir mit ihm einsgemacht und ihm gleich \verbP{geworden sind}\annot{Präsens, Indikativ} in seinem Tod, 

        \spa{2}\konj{so}\annot{Ergebnis} \vb{werden}\annot{Plural, Präsens} wir ihm \konj{auch}\annot{Verbindung} in der Auferstehung gleich \vb{sein}\annot{Infinitiv, Futur I};

        \spa{2}\textbf{6} wir \vb{wissen}\annot{Präsens, Indikativ} ja \konj{dieses}\annot{Bekräftigung}, 

            \spa{3}\konj{dass}\annot{Ursache} unser \person{alter Mensch} \verbP{mitgekreuzigt worden ist}\annot{Präsens, Plural, Indikativ}, 

                \spa{4}\konj{damit}\annot{Begründung} der Leib der Sünde ausser Wirksamkeit \verbP{gesetzt sei}\annot{Präsens, Konjunktiv}, 

                \spa{4}\konj{sodass}\annot{Folge} wir der Sünde nicht mehr \vb{dienen}\annot{Präsens, Indikativ};

            \spa{3}\textbf{7}  \konj{denn}\annot{Auswirkung} wer \vb{gestorben ist}\annot{Perfekt, Indikativ}, der \vb{ist}\annot{Präsens, Indikativ} von der Sünde \vb{freigesprochen}\annot{Präsens, Partizip II}.

\spa{0}\textbf{8} \konj{Wenn}\annot{Bedingung} wir \konj{aber}\annot{Bekräftigung} mit \person{Christus} \verbP{gestorben sind}\annot{Perfekt, Indikativ}, 

    \spa{1}\konj{so} glauben wir, 

        \spa{2}\konj{dass} wir auch mit ihm \vb{leben werden},

        \spa{2}\textbf{9} \konj{da} wir \vb{wissen}, 

            \spa{3}\konj{dass} \person{Christus}, aus den Toten \vb{auferweckt}, nicht mehr \vb{stirbt}; 

            \spa{3}der Tod \vb{herrscht} nicht mehr über ihn.

\spa{2}\textbf{10} \konj{Denn} was er \vb{gestorben ist}, 

    \spa{3}das \vb{ist} er der Sünde \vb{gestorben}, ein für alle Mal; 

    \spa{3}\konj{was} er aber \vb{lebt}, 

    \spa{3}das \vb{lebt} er für \person{Gott}.

\spa{2}\textbf{11} Also auch ihr: \vb{Haltet} euch selbst dafür, 

    \spa{3}\konj{dass} ihr für die Sünde tot \vb{seid}, 

        \spa{4}\konj{aber} für \person{Gott} \vb{lebt} in \person{Christus Jesus, unseren Herrn}!

\spa{0}\textbf{12} \konj{So} \vb{soll} nun die Sünde nicht herrschen in eurem sterblichen Leib, 

    \spa{1}\konj{damit} ihr [der Sünde] nicht durch Begierden [des Leibes] \vb{gehorcht};

    \spa{1}\textbf{13} \vb{gebt} auch nicht eure Glieder der Sünde hin als Werkzeuge der Ungerechtigkeit, 

        \spa{2}\konj{sondern} \vb{gebt} euch selbst \person{Gott} \vb{hin} als solche, 

            \spa{3}die lebendig \vb{geworden sind} aus den \person{Toten}, 

            \spa{3}\konj{und} eure Glieder Gott als Werkzeug der Gerechtigkeit!

    \spa{2}\textbf{14} \konj{Denn} die Sünde \vb{wird} nicht herrschen über euch, 

        \spa{3}\konj{weil} ihr nicht unter dem Gesetz \vb{seid}, 

            \spa{4}\konj{sondern} unter der Gnade.

\spa{0}\textbf{15} \konj{Wie nun}? 

    \spa{1}\vb{sollen} wir \vb{sündigen}, 

    \spa{1}\konj{weil} wir nicht unter dem Gesetz, 

        \spa{2}\konj{sondern}  unter der Gnade \vb{sind}? Das \vb{sei} ferne!

    \spa{1}\textbf{16} \vb{Wisst} ihr nicht: 

    \spa{1}\konj{Wem} ihr euch als \person{Sklaven} \vb{hingebt}, 

        \spa{2}\konj{um} ihm zu \vb{gehorchen}, 

            \spa{3}\konj{dessen} \person{Sklaven} \vb{seid} ihr 

                \spa{4}\konj{und} müsst ihm \vb{gehorchen}, 

                \spa{4}\konj{es} \vb{sei} der Sünde Tode, 

                \spa{4}\konj{oder} dem Gehorsam zur Gerechtigkeit?

\spa{0}\textbf{17} \person{Gott} aber sei Dank, 

    \spa{1}\konj{dass} ihr \person{Sklaven} der Sünde \vb{gewesen}, 

        \spa{2}\konj{nun} aber von Herzen \verbP{gehorsam geworden seid} dem Vorbild der Lehre, 

        \spa{2}das euch \verbP{überliefert worden ist}.

\spa{0}\textbf{18} \konj{Nachdem} ihr aber von der Sünde befreit wurdet, 

    \spa{1}\vb{seid} ihr der Gerechtigkeit dienstbar \vb{geworden}.

\spa{0}\textbf{19} Ich \vb{muss} menschlich davon \vb{reden} 

    \spa{1}\konj{wegen} der Schwachheit eures Fleisches.  

        \spa{2}\konj{Denn so}, wie ihr [einst] eure Glieder in den Dienst der Unreinheit und der Gesetzlosigkeit \vb{gestellt habt} zur Gesetzlosigkeit, 

        \spa{2}\konj{so} \vb{stellt} jetzt eure Glieder in den Dienst der Gerechtigkeit zur Heiligung.

    \spa{1}\textbf{20} \konj{Denn} als ihr \person{Sklaven} der Sünde \vb{wart}, 

    \spa{1}da \vb{wart} ihr frei gegenüber der Gerechtigkeit.

    \spa{1}\textbf{21} \konj{Welche} Frucht \vb{hattet} ihr nun damals von den Dingen, 

        \spa{2}\konj{deren} ihr euch jetzt \vb{schämt}? Ihr Ende ist ja der Tod!

\spa{0}\textbf{22} \konj{Jetzt} aber, 

    \spa{1}\konj{da} ihr von der Sünde frei und \person{Gott} dienstbar \vb{geworden seid}, \vb{habt} ihr 

    \spa{1}\konj{als} eure Frucht die Heiligung, 

    \spa{1}\konj{als} Ende aber ewiges Leben.

\spa{0}\textbf{23} \konj{Denn} der Lohn der Sünde \vb{ist} der Tod; 

    \spa{1}\konj{aber} die Gnadengabe \person{Gottes} \vb{ist} das ewige Leben in \person{Christus Jesus, unserem Herrn}.


\subsubsection{Hinweise}
\begin{itemize}
    \item ??
\end{itemize}

\subsubsection{Beobachtungen}
\begin{itemize}
    \item ??
\end{itemize}

\subsubsection{Aussage / Absicht}
\begin{longtable}{l p{8cm} p{8cm}}
    \hline
    \textbf{Vers} &	\textbf{Aussage} & \textbf{Absicht} \\
    \hline
    \endfirsthead
    \hline
    \textbf{Vers} &	\textbf{Aussage} & \textbf{Absicht} \\
    \hline
    \endhead
    Vers & Aussage & Absicht \\
    \hline
\end{longtable}